\originalchapter{星际介质、恒星形成与冲击波}
本章对应8.901第23--26次课。光谱中的气体既是未来恒星的燃料，也是恒星反馈的接受者。引力、冷却、磁场和角动量共同决定气体能否收缩；一种稳定性阈值并不自动给出完整恒星质量分布。

\originalsection{光电离区以光子预算决定大小}
热星发出的电离光子率$Q_{\rm ion}$照入均匀氢气。若区域内基本电离、稳态且光子没有显著被尘埃吸收，case B把直接到基态的复合光子视为就地再电离，净复合率密度$\alpha_Bn_en_p$。平衡为
\[
 Q_{\rm ion}=\frac{4\pi}3R_S^3\alpha_Bn_en_p,\qquad
 R_S=\left(\frac{3Q_{\rm ion}}{4\pi\alpha_Bn_en_p}\right)^{1/3}.
\]
这是Strömgren半径的静态预算，不是前沿任何时刻都立即达到它。复合时间$t_{\rm rec}=1/(\alpha_Bn_e)$比较于源寿命和膨胀时间才可判断稳态是否合理。非均匀气体的复合取$\langle n_en_p\rangle$，不能把平均密度平方直接替代；膨胀和流体响应又会改变体积。

中性原子氢21 cm超精细跃迁可示踪HI，自旋温度连接上下能级布居，吸收与发射取决于相对背景亮温。光学薄、激发受控时积分线亮度可转为原子氢柱密度，光学厚或冷前景会改变推断。冷H$_2$的对称性使低温易激发的强偶极线缺少，常用CO或尘埃间接示踪；其转换系数还受金属、激发、光学深度和环境影响，不是无条件的通用质量计。

\originalsection{热平衡与热稳定是不同问题}
设体积净冷却$\mathcal L=n^2\Lambda(T)-n\Gamma(T)$，正号表示失热；平衡$\mathcal L=0$只给一个候选温度。若声波调整快，可近等压，此时$n=P/(k_BT)$。每粒子的净损失$\ell=\mathcal L/n=n\Lambda-\Gamma$，焓方程$c_P\dot T=-\ell$。在平衡附近
\[
 c_P\delta\dot T=-\left(\frac{\partial\ell}{\partial T}\right)_P\delta T.
\]
所以$\ell'_P>0$才稳定：温度稍高就冷却得更多。若$\Gamma$近常数，判据化为$T\Lambda'(T)-\Lambda(T)>0$。等密度扰动用另一导数，不能把一个图上的相交点都当成稳定气相。真实气体还有传导、化学、辐射和时间依赖。

冷却时间大致为热能密度除净损失；光学薄、理想单原子时可用$(3/2)n_{\rm tot}k_BT/(n_en_i\Lambda)$，但有外加热时分母应为净值。PS6的宇宙暗晕冷却题特别区分密度、电子离子数及能量如何分给少量重子，这些条件同样影响星际气体的时钟。

\originalsection{Jeans阈值只是收缩的入口}
第\ref{ch:compact}章已由连续、Euler、Poisson推$\omega^2=c_s^2k^2-4\pi G\rho_0$。本卷取$\lambda_J=2\pi c_s/\sqrt{4\pi G\rho_0}$、球半径$R_J=\lambda_J/2$，所以$M_J=\pi^{5/2}c_s^3/(6G^{3/2}\rho_0^{1/2})$。符号$c_s$是选定扰动的等温或绝热响应，必须和热交换时间一致。

等温坍缩保持$c_s$近常数，$M_J\propto\rho^{-1/2}$下降，可能进一步碎裂。若气体变光学厚、不能快速失热，温度上升改变趋势。线性Jeans模型采用背景减引力的无限均匀理想化；湍流、有限边界、磁场和旋转需要额外条件，不能仅用一个$M_J$推出每个真实云的最终恒星质量。

磁通近冻结时$B\propto R^{-2}$，磁能量$\sim B^2R^3\propto R^{-1}$，与引力$|W|\sim GM^2/R$同尺度。故质量与磁通比$M/\Phi_B$而非单看$B$决定磁支持重要性，临界系数取决于几何和边界。非理想磁扩散可改变磁通保存。

比角动量$j$近守恒的质点有离心平衡尺度$r_{\rm circ}=j^2/(GM)$。PS10.5的第一次径向转向对近零总能轨道给$j^2/(2GM)$，它与圆化半径不同；前者保留径向能量，后者需耗散使轨道近圆。形成盘之后还需角动量向外输运才能继续吸积。

棕矮星在收缩加热足以维持长期氢燃烧前受电子简并影响，不能简单用理想气体温度随$M/R$无止境升高。初始质量函数则受碎裂、吸积、反馈和统计选择影响；Jeans量级或单个多方模型并未推导出其完整形状。

\originalsection{冲击波必须同时满足三个通量守恒}
在平面、稳态、无额外辐射损失的冲击静止系，上下游法向速度$u_1,u_2$：
\[
 \rho_1u_1=\rho_2u_2,\quad
 P_1+\rho_1u_1^2=P_2+\rho_2u_2^2,\quad
 \frac{u_1^2}{2}+\frac{\gamma P_1}{(\gamma-1)\rho_1}
 =\frac{u_2^2}{2}+\frac{\gamma P_2}{(\gamma-1)\rho_2}.
\]
强冲击$P_1\ll\rho_1u_1^2$时，记压缩比$x=\rho_2/\rho_1$，连续给$u_2=u_1/x$，动量给$P_2=\rho_1u_1^2(1-1/x)$。代能量式，除共同$u_1^2/2$，可得非平凡根$x=(\gamma+1)/(\gamma-1)$；另一个$x=1$是不发生冲击。单原子气体$\gamma=5/3$给$x=4$，
\[
 P_2=\frac34\rho_1u_1^2,\qquad
 k_BT_2=\frac{3}{16}\mu m_u u_1^2.
\]
高速运动变成随机热运动，温度随速度平方；电子与离子不一定立即同温，强磁场、相对论粒子或显著冷却都会改简单跳变。

\originalsection{超新星遗迹的阶段用预算区分}
最早扫过质量小于抛出物质量时可近自由膨胀。扫过介质占主导且能量$E$近守恒时，唯一长度由$E,\rho_0,t$组成，
\[
 R(t)=\xi\left(\frac{Et^2}{\rho_0}\right)^{1/5},\qquad
 \dot R=\frac25R/t.
\]
$\xi$须求相似流体解，量纲分析不决定它。冲击加热、辐射和宇宙线影响总预算；若后期能量大量辐射，不能继续用同一Sedov尺度。近守恒壳动量$\rho_0R^3\dot R$的极限给$R\propto t^{1/4}$，与能量守恒$t^{2/5}$不同。

遗迹推动湍流、加热介质、输送元素，也可能压缩冷云。这类反馈将单星演化接回下一代恒星形成；其净促进或抑制效果取决于冷却、几何和云状态，不由“有冲击”三个字自动判断。
